\documentclass[10pt,a4paper]{amsart}
\usepackage[margin=19mm]{geometry}
\usepackage{amsmath,amssymb,mathtools}
\usepackage[scr=rsfs,cal=euler]{mathalfa}
\usepackage{xfrac}
\usepackage[unicode,hidelinks,bookmarks=false]{hyperref}
\newcommand{\ie}{i.e.\ }
\newcommand{\cf}{cf.\ }
\newcommand{\resp}{respectively\ }
\newcommand{\Subsection}{\subsection}
\providecommand{\nobreakdash}{\nobreak\hbox{-}}
\setlength{\emergencystretch}{2em}
\multlinegap0pt
\usepackage{mdframed}
\usepackage{mdframed}
\usepackage{mdframed}
\usepackage{mdframed}
\usepackage{amssymb}
\usepackage{mdframed}
\usepackage{algorithm}\usepackage{algpseudocode}
\usepackage{dsfont}
\usepackage[shortlabels]{enumitem}
\usepackage{mathtools}
\usepackage{cancel}
\newtheorem{thm}{Theorem}
\newcommand{\N}{\mathbb{N}}
\newcommand{\R}{\mathbb{R}}
\title{Lecture notes: Stochastic gradient methods}
\author{Simon Weissmann}
\date{Source: arXiv:2606.03953v1 (CC BY 4.0)}
\begin{document}
\maketitle

\noindent\emph{Source and licence.} Lecture notes: Stochastic gradient methods. Version arXiv:2606.03953v1 (CC BY 4.0), distributed by arXiv under the Creative Commons Attribution 4.0 International licence, http://creativecommons.org/licenses/by/4.0/, as recorded in the arXiv licence metadata for this exact version and shown on the abstract page https://arxiv.org/abs/2606.03953v1. Full licence notice: \texttt{licenses/arxiv-2606.03953-CC-BY-4.0.txt}.\par\smallskip

\noindent\textit{Editorial scope.} Counted unit: the article's thm thm (source0.tex lines 1936--1939). The article's complete original proof of it is delivered with it (source0.tex lines 1941--2010, one \texttt{proof} environment of 70 lines). No mathematics is written, repaired or re-derived: the statement and the proof are the article's own lines, in the article's own notation, and no retained line is substituted or re-flowed.  Retained uncounted as the source-local context the proof needs: lines 1890--1896.  Nothing was omitted from any retained span, so the package carries no omission record.  No retained line contains a citation, so the wrapper carries no bibliography and no bibliography entry.  No retained line contains a figure environment, an image-inclusion command or drawing code, so no image is delivered and the package declares no asset dependency.  Editorial formatting only, in addition: an AMS \texttt{amsart} preamble that restates the article's own declarations and macros used by the retained lines, namely the environments \texttt{thm} on separate counters, together with the standard packages those lines need.  The retained environments print in this document as Theorem 1 in order of appearance, whereas the article prints the same items with the numbering of its own full document.  The complete native source of the article as distributed in the arXiv e-print for this exact version is bundled unchanged as \texttt{sources/202014/source0.tex}.
\begin{equation}\label{eq:NAM_strongconvex}
\begin{split}
x_k &= \frac{\tau}{1+\tau} z_k + \frac{1}{1+\tau} y_k\\
y_{k+1} &= x_k - \frac{1}{L}\nabla f(x_k)\\
z_{k+1} &= z_k +\tau(x_k-z_k)-\frac{\tau}{\mu}\nabla f(x_k)
\end{split}\, 
\end{equation}

\begin{thm}[NAM for strongly convex and smooth objective function] 
Let $f:\R^d\to\R$ be $\mu$-strongly convex and $L$-smooth with $L>\mu$, and let $x_\ast\in\R^d$ be the corresponding unique global minimum of $f$. Moreover, let $(x_k,y_k,z_k)_{k\in\N}$ be generated by \eqref{eq:NAM_strongconvex} with $\tau = \sqrt{\frac{\mu}{L}}\in(0,1)$ and initialized by $(y_0,z_0)\in\R^d\times\R^d$. Then {NAM} converges linearly in the sense that
\[e_k := f(y_k)-f(x_\ast) + \frac{\mu}{2}\|z_k-x_\ast\|^2 \le \left(1-\sqrt{\frac{\mu}{L}}\right)^k \left(f(y_0)-f(x_\ast) +\frac{\mu}{2}\|z_0-x_\ast\|^2\right). \]
\end{thm}

\begin{proof}
We define $e_k := f(y_k)-f(x_\ast) + \frac{\mu}{2}\|z_k-x_\ast\|^2$ and aim to prove 
\[e_{k+1}\le (1-\tau) e_k. \]
Firstly, applying $L$-smoothness of $f$ gives
\begin{equation}\label{eq:NAM_strongconvex_aux1}
f(y_{k+1}) \le f(x_k) + \langle \nabla f(x_x), y_{k+1}-x_k\rangle + \frac{L}{2}\|x_k-y_{k+1}\|^2 = f(x_k) -\frac{1}{2L}\|\nabla f(x_k)\|^2.
\end{equation}
Next, we consider the update of $e_k^{(1)} = \|z_k-x_\ast\|^2$:
\begin{align*}
e_{k+1}^{(1)} &= \|z_k +\tau(x_k-z_k) -\frac{\tau}{\mu}\nabla f(x_k) - x_\ast\|^2\\
&= \|z_k-x_\ast\|^2 + 2\langle \tau (x_k-z_k)-\frac{\tau}{\mu}\nabla f(x_k), z_k - x_\ast\rangle + \|\underbrace{\tau(x_k-z_k) - \frac{\tau}{\mu}\nabla f(x_k)}_{=z_{k+1}-z_k} \|^2\\
&= e_k^{(1)} + 2\tau \langle x_k-z_k,\underbrace{z_k-x_\ast}_{=z_k-x_k+x_k-x_\ast}\rangle - 2\frac{\tau}{\mu} \langle \nabla f(x_k),\underbrace{z_k-x_\ast}_{=z_k-x_k+x_k-x_\ast}\rangle +\|z_{k+1}-z_k\|^2\\
&= e_k^{(1)}+ \|z_{k+1}-z_k\|^2 + 2\tau \langle x_k-z_k,z_k-x_k\rangle - 2\frac{\tau}{\mu}\langle \nabla f(x_k),z_k-x_k\rangle \\
&\quad+2\tau \langle x_k-z_k,x_k-x_\ast\rangle -2\frac{\tau}{\mu}\langle \nabla f(x_k), x_k-x_\ast\rangle 
\end{align*}
Recall that by strong convexity it holds true that
\[f(x_\ast)-f(x_k)\ge \langle x_\ast-x_k,\nabla f(x_k)\rangle + \frac{\mu}2 \|x_k-x_\ast\|^2, \]
such that
\begin{align*}
e_{k+1}^{(1)} &\le e_k^{(1)}+ \|z_{k+1}-z_k\|^2 + 2\tau \langle x_k-z_k,z_k-x_k\rangle - 2\frac{\tau}{\mu}\langle \nabla f(x_k),z_k-x_k\rangle \\
&\quad + 2\frac{\tau}{\mu} (f(x_\ast)-f(x_k)-\frac{\mu}2\|x_k-x_\ast\|^2) + 2\tau \langle x_k-z_k,x_k-x_\ast\rangle.
\end{align*}
We observe that
\[\tau(x_k-z_k) = \tau x_k -((1+\tau) x_k -y_k) = y_k-x_k \]
and obtain
\begin{align*} e_{k+1}^{(1)} &\le e_k^{(1)}+ \|z_{k+1}-z_k\|^2-2\frac{\tau}{\mu} (f(x_k)-f(x_\ast))+ \frac{2}{\mu}\langle \nabla f(x_k),y_k-x_k\rangle\\
&\quad+ 2\tau\langle x_k-z_k,x_k-x_\ast\rangle -2\tau\|x_k-z_k\|^2 -\tau\|x_k-x_\ast\|^2.
\end{align*}
Note that
\begin{align*}2\tau\langle x_k-z_k,x_k-x_\ast\rangle -2\tau\|x_k-z_k\|^2 -\tau\|x_k-x_\ast\|^2 &= -\tau \|x_k-x_\ast-(x_k-z_k)\|^2 - \tau \|x_k-z_k\|^2\\ &= -\tau e_k^{(1)} -\tau\|x_k-z_k\|^2,
\end{align*}
which yields
\begin{align*} 
e_{k+1}^{(1)} &\le (1-\tau) e_k^{(1)}+ \|z_{k+1}-z_k\|^2-2\frac{\tau}{\mu} (f(x_k)-f(x_\ast))+ \frac{2}{\mu}\langle \nabla f(x_k),y_k-x_k\rangle-\tau\|x_k-z_k\|^2.
\end{align*}
Finally, with $e_k^{(2)}=f(y_k)-f(x_\ast)$ we consider the evolution of $e_k$ through
\begin{align*}
e_{k+1} &= \frac{\mu}{2}e_{k+1}^{(1)}+ e_{k+1}^{(2)}\le (1-\tau)\frac{\mu}{2}e_k^{(1)}+\frac{\mu}{2}\|z_{k+1}-z_k\|^2-\tau (f(x_k)-f(x_\ast))+\langle \nabla f(x_k),y_k-x_k\rangle\\ &\quad-\tau\frac{\mu}{2}\|x_k-z_k\|^2 + f(y_{k+1})-f(x_\ast)\\
&\le  (1-\tau)\frac{\mu}{2}e_k^{(1)}+\frac{\mu}{2}\|z_{k+1}-z_k\|^2-\tau (f(x_k)-f(x_\ast))+\langle \nabla f(x_k),y_k-x_k\rangle\\ &\quad-\tau\frac{\mu}{2}\|x_k-z_k\|^2 + f(x_{k})-f(x_\ast)-\frac{1}{2L}\|\nabla f(x_k)\|^2+\{(1-\tau)e_k^{(2)}-(1-\tau)e_k^{(2)}\}  \}\\
&= (1-\tau) e_k -(1-\tau) (f(y_k)-f(x_\ast))-\tau (f(x_k)-f(x_\ast)) + f(x_k)-f(x_\ast)-\frac{1}{2L}\|\nabla f(x_k)\|^2\\
&\quad +\frac{\mu}{2}\|z_{k+1}-z_k\|^2-\tau\frac{\mu}{2}\|x_k-z_k\|^2+\langle \nabla f(x_k),y_k-x_k\rangle\\
&=:  (1-\tau) e_k  + \mathcal R,
\end{align*}
where we have used \eqref{eq:NAM_strongconvex_aux1} in the first inequality and added a zero $(1-\tau) (f(y_k)-f(x_\ast))-(1-\tau) (f(y_k)-f(x_\ast))=0$. It is left to prove that $\mathcal R\le 0$. First note, that $\mathcal R$ simplifies to
\[ \mathcal R = (1-\tau)(f(x_k)-f(y_k)) - \frac1{2L}\|\nabla f(x_k)\|^2 +\frac{\mu}{2}\|z_{k+1}-z_k\|^2-\tau\frac{\mu}{2}\|x_k-z_k\|^2+(1-\tau+\tau)\langle \nabla f(x_k),y_k-x_k\rangle.\]
We apply strong convexity in the form of 
\[\langle \nabla f(x_k),y_k-x_k\rangle \le f(y_k) - f(x_k) -\frac{\mu}{2}\|y_k-x_k\|^2 \]
such that
\begin{align*}
\mathcal R &\le (1-\tau)(f(x_k)-f(y_k))+(1-\tau)(f(y_k)-f(x_k))-\frac{\mu}{2}\|y_k-z_k\|^2\\
&\quad -\frac{1}{2L}\|\nabla f(x_k)\|^2 +\frac{\mu}{2}\|z_{k+1}-z_k\|^2-\tau\frac{\mu}{2}\|x_k-z_k\|^2+\tau\langle \nabla f(x_k),y_k-x_k\rangle \\
&=\tau\langle \nabla f(x_k),y_k-x_k\rangle-(1-\tau)\frac{\mu}{2}\|y_k-x_k\|^2-\frac{1}{2L}\|\nabla f(x_k)\|^2\\
&\quad+\frac{\mu}{2}\|z_{k+1}-z_k\|^2-\tau\frac{\mu}{2}\|x_k-z_k\|^2.
\end{align*}
Recall that we have used $z_{k+1}-z_k = \tau(x_k-z_k)-\frac{\tau}{\mu}\nabla f(x_k)$, which with $\tau(x_k-z_k) = y_k-x_k$ gives
\[\frac{\mu}{2} \|z_{k+1}-z_k\|^2 = \frac{\mu}2\|y_k-x_k\|^2 -\tau \langle \nabla f(x_k),y_k-x_k\rangle + \frac{\tau^2}{2\mu}\|\nabla f(x_k)\|^2, \]
and therefore, we obtain
\begin{align*}
\mathcal R &\le -(1-\tau)\frac{\mu}{2} \|y_k-x_k\|^2 +\frac{\mu}2\|y_k-x_k\|^2 + \tau\langle \nabla f(x_k),y_k-x_k\rangle -\tau \langle \nabla f(x_k),y_k-x_k\rangle\\
&\quad +\frac{\tau^2}{2\mu}\|\nabla f(x_k)\|^2 - \frac{1}{2L}\|\nabla f(x_k)\|^2 - \tau\frac{\mu}{2}\frac{\tau}{\tau}\|x_k-z_k\|^2\\
&=\tau\frac{\mu}{2}\|y_k-x_k\|^2 - \frac{\mu}{2\tau}\|\tau(x_k-z_k)\|^2 + \left(\frac{\tau^2}{2\mu}-\frac{1}{2L} \right)\|\nabla f(x_k)\|^2\\
&= \left(\frac{\tau\mu}{2}-\frac{\mu}{2\tau}\right) \|y_k-x_k\|^2 +\left(\frac{\tau^2}{2\mu}-\frac{1}{2L} \right)\|\nabla f(x_k)\|^2,
\end{align*}
where we have used again that $\tau(x_k-z_k) = y_k-x_k$. With the choice $\tau=\sqrt{\frac{\mu}{L}}\in(0,1)$ we observe that
\[ \frac{\tau^2}{2\mu}-\frac{1}{2L} = \frac{1}{2L}-\frac{1}{2L} = 0\]
and
\[ \frac{\tau\mu}{2}-\frac{\mu}{2\tau} = \frac{\mu}2(\tau-\frac1\tau) \le 0.\]
This proves that $\mathcal R\le 0$ and the assertion follows:
\[e_{k+1} \le (1-\tau)e_k = \left(1-\sqrt{\frac{\mu}{L}}\right)e_k. \]
\end{proof}

\end{document}
